\newcommand{\scalefig}{0.5}


% this command takes the language as an argument
% the default is Bash
% it sets the other format to be as described in the function

% List of possible programming languages to use
% ABAP, ACSL, Ada, Algol, Ant, Assembler, Awk, bash, Basic, C#, C++, C, Caml, Clean, Cobol, Comal, csh, Delphi, Eiffel, Elan, erlang, Euphoria, Fortran4, GCL, Go (golang), Gnuplot, Haskell, HTML, IDL, inform, Java, JVMIS, ksh, Lisp, Logo, Lua, make, Mathematica1, Matlab, Mercury, MetaPost, Miranda, Mizar, ML, Modelica, Modula-2, MuPAD, NASTRAN, Oberon-2, Objective C , OCL, Octave, Oz, Pascal, Perl, PHP, PL/I, Plasm, POV, Prolog, Promela, Python, R, Reduce, Rexx, RSL, Ruby, S, SAS, Scilab, sh, SHELXL, Simula, SQL, tcl, TeX, VBScript, Verilog, VHDL, VRML, XML, XSLT.
\definecolor{dkgreen}{rgb}{0,0.6,0}
\newcommand{\setCustomListing}[1]{

    
    \lstset{
      language={#1},
      basicstyle={\small},
      identifierstyle={\small},
      commentstyle={\small\itshape\color{dkgreen}},
      keywordstyle={\small\bfseries\color{blue}},
      ndkeywordstyle={\small},
      stringstyle={\small\ttfamily},
      frame={tb},
      breaklines=true,
      columns=[l]{fullflexible},
      numbers=left,
      xrightmargin=0em,
      xleftmargin=3em,
      numberstyle={\scriptsize},
      stepnumber=1,
      numbersep=1em,
      lineskip=-0.5ex,
    }



}

%% setCustomListing Command End

%% ArithV Language Definition Begin %%
\lstdefinelanguage{ArithV}{
  keywords=[1]{input, output, variable, TestBench, Module},  % Type keywords
  keywords=[2]{Add, Sub, Mul, Shl, Shr, And, Or, Xor, Not, Concat, Slice},  % Operation keywords
  keywords=[3]{is_signed, op_mode},  % Attribute keywords
  sensitive=true,
  comment=[l]{\#},
  morestring=[b]",
  morestring=[b]',
}

\definecolor{arithvpurple}{rgb}{0.5,0,0.5}
\definecolor{arithvred}{rgb}{0.6,0,0}
\definecolor{arithvorange}{rgb}{0.8,0.4,0}

\lstdefinestyle{arithvstyle}{
  language=ArithV,
  basicstyle={\small\ttfamily},
  identifierstyle={\small\ttfamily},
  commentstyle={\small\itshape\color{dkgreen}},
  keywordstyle=[1]{\small\bfseries\color{blue}},      % input, output, variable
  keywordstyle=[2]{\small\bfseries\color{arithvpurple}},  % Operations
  keywordstyle=[3]{\small\color{arithvorange}},       % Attributes
  stringstyle={\small\ttfamily\color{arithvred}},
  frame={tb},
  breaklines=true,
  columns=[l]{fullflexible},
  numbers=left,
  xrightmargin=0em,
  xleftmargin=3em,
  numberstyle={\scriptsize\color{gray}},
  stepnumber=1,
  numbersep=1em,
  lineskip=-0.5ex,
  showstringspaces=false,
  literate={×}{x}1 {→}{->}2 {←}{<-}2,
  inputencoding=utf8,
  extendedchars=true,
}
%% ArithV Language Definition End %%

%% Minted ArithV Support Begin %%
% Requires: -shell-escape flag and arithvlexer.py installed
% To install: /opt/homebrew/opt/python@3.10/bin/python3.10 -m pip install -e . --user
% See README_MINTED_SETUP.md for full setup instructions
\usepackage{minted}

% Use colorful style for better syntax highlighting
\usemintedstyle{colorful}

% Define colors for Pygments tokens to match lstlisting arithvstyle
\makeatletter
\expandafter\def\csname PYG@tok@kt\endcsname{\def\PYG@tc##1{\textcolor[rgb]{0,0,1}{##1}}}        % Keyword.Type -> blue
\expandafter\def\csname PYG@tok@kr\endcsname{\def\PYG@tc##1{\textcolor[rgb]{0.5,0,0.5}{##1}}}   % Keyword.Reserved -> purple
\expandafter\def\csname PYG@tok@na\endcsname{\def\PYG@tc##1{\textcolor[rgb]{0.8,0.4,0}{##1}}}   % Name.Attribute -> orange
\expandafter\def\csname PYG@tok@s\endcsname{\def\PYG@tc##1{\textcolor[rgb]{0.6,0,0}{##1}}}      % String -> red
\expandafter\def\csname PYG@tok@c\endcsname{\def\PYG@tc##1{\textcolor[rgb]{0,0.6,0}{\textit{##1}}}} % Comment -> green italic
\makeatother

\newmintedfile[arithvfile]{arithv}{%
  breaklines=true,
  fontsize=\small,
  fontfamily=tt,
  frame=lines,
  linenos=true,
  numbersep=1em,
  xleftmargin=3em,
  autogobble=true,
  breaksymbolleft={},
  breaksymbolright={},
}

\newminted{arithv}{%
  breaklines=true,
  fontsize=\small,
  fontfamily=tt,
  frame=lines,
  linenos=true,
  numbersep=1em,
  xleftmargin=3em,
  autogobble=true,
  breaksymbolleft={},
  breaksymbolright={},
}

% Alias for ArithV minted environments (alternative to lstinputlisting)
\newcommand{\arithvinputlisting}[2][]{%
  \inputminted[#1]{arithv}{#2}%
}
%% Minted ArithV Support End %%





% \usepackage{todonotes}  % Already loaded in s0_preamble.tex

% Simple colored todo in margin (supports complex content like itemize)
\newcommand{\reptodo}[1]{\marginpar{\raggedright\textcolor{red}{\begin{minipage}{\marginparwidth}\small #1\end{minipage}}}}